\documentclass[12pt,a4paper]{exam}
\usepackage[utf8]{inputenc}
\usepackage{amsmath,amssymb,amsfonts}
\usepackage{geometry}
\usepackage{enumitem}
\usepackage{graphicx}
\usepackage{hyperref}
\usepackage{xcolor}
\geometry{a4paper, top=2cm, bottom=2cm, left=2.5cm, right=2.5cm}
\hypersetup{colorlinks=true, linkcolor=blue, urlcolor=blue}
\renewcommand{\thequestion}{\textbf{\arabic{question}}}
\renewcommand{\questionlabel}{\thequestion.}
\pointname{ marks}
\pointsdroppedatright
\bracketedpoints
\setlength{\parindent}{0pt}
\setlength{\emergencystretch}{3em}
\allowdisplaybreaks
\newsavebox{\schmathbox}
\newcommand{\fitmath}[1]{%
  \sbox{\schmathbox}{$\displaystyle #1$}%
  \ifdim\wd\schmathbox>\linewidth
    \resizebox{\linewidth}{!}{\usebox{\schmathbox}}%
  \else
    \usebox{\schmathbox}%
  \fi}

\begin{document}

\thispagestyle{empty}
\begin{center}
  \textbf{\LARGE Diag} \\[0.3cm]
  \textbf{Time Allowed: 3 Hours} \hfill \textbf{Maximum Marks: 80} \\[0.3cm]
  \rule{\textwidth}{0.5pt}
\end{center}

\vspace{0.3cm}

\noindent\textbf{General Instructions:}
\begin{enumerate}
  \item {\normalfont\normalcolor This question paper contains 36 questions. All questions are compulsory.}
  \item {\normalfont\normalcolor Questions are grouped into sections A through E.}
  \item {\normalfont\normalcolor Use of calculators is NOT permitted.}
\end{enumerate}
\bigskip
\noindent\textbf{\large Section A: Multiple Choice \& Assertion-Reason (1 mark each)}

\begin{questions}
\question[1] If $x = 2$ is a root of $x^2 + kx + 4 = 0$, the value of $k$ is
\begin{enumerate}[label=(\Alph*)]
  \item (A) 4
  \item (B) 2
  \item (C) -4
  \item (D) -2
\end{enumerate}
\vspace{0.15cm}

\question[1] If the HCF of two numbers $a$ and $b$ is $12$ and their product is $1800$, then their LCM is:
\begin{enumerate}[label=(\Alph*)]
  \item (A) 120
  \item (B) 140
  \item (C) 150
  \item (D) 160
\end{enumerate}
\vspace{0.15cm}

\question[1] The decimal expansion of the rational number $\frac{43}{2^4 \times 5^3}$ will terminate after how many places of decimal?
\begin{enumerate}[label=(\Alph*)]
  \item (A) 3
  \item (B) 4
  \item (C) 2
  \item (D) 1
\end{enumerate}
\vspace{0.15cm}

\question[1] If $p$ and $q$ are two prime numbers, then the HCF of $p^3q^2$ and $p^2q$ is:
\begin{enumerate}[label=(\Alph*)]
  \item (A) p\textasciicircum{}2q
  \item (B) p\textasciicircum{}3q\textasciicircum{}2
  \item (C) pq
  \item (D) p\textasciicircum{}2q\textasciicircum{}2
\end{enumerate}
\vspace{0.15cm}

\question[1] If two positive integers $a$ and $b$ are written as $a = x^3y^2$ and $b = xy^3$, where $x, y$ are prime numbers, then $LCM(a, b)$ is:
\begin{enumerate}[label=(\Alph*)]
  \item (A) xy
  \item (B) x\textasciicircum{}2y\textasciicircum{}2
  \item (C) x\textasciicircum{}3y\textasciicircum{}3
  \item (D) x\textasciicircum{}4y\textasciicircum{}5
\end{enumerate}
\vspace{0.15cm}

\question[1] The smallest number which is divisible by both $306$ and $657$ is:
\begin{enumerate}[label=(\Alph*)]
  \item (A) 12338
  \item (B) 22338
  \item (C) 21338
  \item (D) 23338
\end{enumerate}
\vspace{0.15cm}

\question[1] The discriminant of the quadratic equation $x^2 - 5x + 6 = 0$ is:
\begin{enumerate}[label=(\Alph*)]
  \item (A) 1
  \item (B) -1
  \item (C) 49
  \item (D) 0
\end{enumerate}
\vspace{0.15cm}

\question[1] The discriminant of the quadratic equation $2x^2 - 3x + 1 = 0$ is:
\begin{enumerate}[label=(\Alph*)]
  \item (A) 1
  \item (B) -1
  \item (C) 9
  \item (D) -7
\end{enumerate}
\vspace{0.15cm}

\question[1] The discriminant of the quadratic equation $2x^2 - 3x + 5 = 0$ is:
\begin{enumerate}[label=(\Alph*)]
  \item (A) $31$
  \item (B) $-31$
  \item (C) $41$
  \item (D) $-41$
\end{enumerate}
\vspace{0.15cm}

\question[1] The roots of the quadratic equation $x^2 - 7x + 10 = 0$ are:
\begin{enumerate}[label=(\Alph*)]
  \item (A) 2, 5
  \item (B) -2, -5
  \item (C) 2, -5
  \item (D) -2, 5
\end{enumerate}
\vspace{0.15cm}

\question[1] Which of the following is the discriminant of the quadratic equation $2x^2 - 4x + 3 = 0$?
\begin{enumerate}[label=(\Alph*)]
  \item (A) $\sqrt{40}$
  \item (B) $4$
  \item (C) $16 - 24$
  \item (D) $0$
\end{enumerate}
\vspace{0.15cm}

\question[1] Assertion (A): The quadratic equation $2x^2 - 5x + 3 = 0$ has two distinct real roots.
Reason (R): For a quadratic equation $ax^2 + bx + c = 0$, if $b^2 - 4ac > 0$, then the equation has two distinct real roots.
\begin{enumerate}[label=(\Alph*)]
  \item (A) Both Assertion (A) and Reason (R) are true and Reason (R) is the correct explanation of Assertion (A)
  \item (B) Both Assertion (A) and Reason (R) are true but Reason (R) is NOT the correct explanation of Assertion (A)
  \item (C) Assertion (A) is true but Reason (R) is false
  \item (D) Assertion (A) is false but Reason (R) is true
\end{enumerate}
\vspace{0.15cm}

\question[1] Assertion (A): The quadratic equation $x^2 - 3x + 5 = 0$ has no real roots.
Reason (R): For a quadratic equation $ax^2+bx+c=0$, if discriminant $D = b^2 - 4ac < 0$, then the equation has no real roots.
\begin{enumerate}[label=(\Alph*)]
  \item (A) Both Assertion (A) and Reason (R) are true and Reason (R) is the correct explanation of Assertion (A)
  \item (B) Both Assertion (A) and Reason (R) are true but Reason (R) is NOT the correct explanation of Assertion (A)
  \item (C) Assertion (A) is true but Reason (R) is false
  \item (D) Assertion (A) is false but Reason (R) is true
\end{enumerate}
\vspace{0.15cm}

\end{questions}
\bigskip
\noindent\textbf{\large Section B: Very Short Answer (2 marks each)}

\begin{questions}
\question[1] Find the discriminant of the quadratic equation $2x^2 - 4x + 3 = 0$.
\vspace{0.15cm}

\question[1] Find the discriminant of the quadratic equation $2x^2 - 3x + 1 = 0$.
\vspace{0.15cm}

\question[1] Find the discriminant of the quadratic equation $2x^2 - 4x + 3 = 0$.
\vspace{0.15cm}

\question[2] Find the discriminant of the quadratic equation $2x^2 - 4x + 3 = 0$. Hence, determine the nature of its roots.
\vspace{0.15cm}

\question[2] Find the nature of the roots of the quadratic equation $2x^2 - 4x + 3 = 0$.
\vspace{0.15cm}

\question[2] Find the nature of the roots of the quadratic equation $2x^2 - 4x + 3 = 0$.
\vspace{0.15cm}

\question[2] Find the discriminant of the quadratic equation $3x^2 - 5x + 2 = 0$ and hence determine the nature of its roots.
\vspace{0.15cm}

\question[2] Find the discriminant of the quadratic equation $3x^2 - 6x + 2 = 0$ and hence determine the nature of its roots.
\vspace{0.15cm}

\end{questions}
\bigskip
\noindent\textbf{\large Section C: Short Answer (3 marks each)}

\begin{questions}
\question[3] A train travels 360 km at a uniform speed. If the speed had been 5 km/h more, it would have taken 1 hour less for the same journey. Find the speed of the train.
\vspace{0.15cm}

\question[3] The speed of a boat in still water is 8 km/h. It takes the same time to travel 15 km upstream as it takes to travel 25 km downstream. Find the speed of the stream.
\vspace{0.15cm}

\question[3] The sum of the squares of two consecutive positive integers is 365. Find the integers.
\vspace{0.15cm}

\question[3] The sum of two numbers is 15. If the sum of their reciprocals is $\frac{3}{10}$, find the numbers.
\vspace{0.15cm}

\question[3] A train travels 360 km at a uniform speed. If the speed had been 5 km/h more, it would have taken 1 hour less for the same journey. Find the speed of the train.
\vspace{0.15cm}

\end{questions}
\bigskip
\noindent\textbf{\large Section D: Long Answer (4-5 marks each)}

\begin{questions}
\question[4] If the polynomial $p(x) = x^4 - 6x^3 + 16x^2 - 25x + 10$ is divided by $g(x) = x^2 - 2x + k$, the remainder is found to be $x + a$. What is the value of $k$?
\vspace{0.15cm}

\question[1] When $f(x) = x^3 - 3x^2 + 5x - 3$ is divided by $g(x) = x^2 - 2$, the quotient $q(x)$ is:
\vspace{0.15cm}

\question[1] If $g(x)$ is the divisor of $p(x) = 2x^4 - 3x^3 - 3x^2 + 6x - 2$ and the quotient $q(x) = 2x^2 - 3x + 1$ with $r(x) = 0$, what is the divisor $g(x)$?
\vspace{0.15cm}

\question[1] For a polynomial $p(x)$, if $p(x) = (x^2 - 4)q(x) + (2x + 1)$, what is the remainder when $p(x)$ is divided by $(x - 2)$?
\vspace{0.15cm}

\question[1] A polynomial $p(x)$ of degree $3$ is divided by $(x^2 - 1)$ resulting in a quotient $(x - 1)$ and a remainder $(2x + 3)$. Find the polynomial $p(x)$.
\vspace{0.15cm}

\question[5] A train travels a distance of 480 km at a uniform speed. If the speed had been 8 km/h less, then it would have taken 3 hours more to cover the same distance. Find the original speed of the train. Also, find the time taken by the train to cover the distance at the original speed.
\vspace{0.15cm}

\question[5] A train travels a distance of 480 km at a uniform speed. If the speed had been 8 km/h less, then it would have taken 3 hours more to cover the same distance. Find the original speed of the train. Also, find the time taken at the original speed.
\vspace{0.15cm}

\question[5] A train travels a distance of 480 km at a uniform speed. If the speed had been 8 km/h less, then it would have taken 3 hours more to cover the same distance. Find the original speed of the train. Also, find the time taken at the original speed.
\vspace{0.15cm}

\question[5] A motorboat whose speed in still water is $18$ km/h takes $1$ hour more to go $24$ km upstream than to return downstream to the same spot. Find the speed of the stream.
\vspace{0.15cm}

\question[5] A motorboat whose speed in still water is 18 km/h takes 1 hour more to go 48 km upstream than to return to the same point downstream. Find the speed of the stream.
\vspace{0.15cm}

\end{questions}
\bigskip

\end{document}
